% Original Irodov Part V §5.1 36--45, printed pp.193--194.
% Ray diagram 5.9 used for three-thin-lens system.
\opticsproblem{36}{193}
点光源与屏幕间固定距离 $l=90$cm，
一只薄会聚透镜在两个不同位置都可投得清晰实像。
求透镜焦距：
(a) 两个透镜位置间隔 $\Delta l=30$cm；
(b) 两位置得到的像横向大小之比 $\eta=4.0$。
\exsolution
两共轭透镜位置互换物距、像距：
$s+s'=l$，$|s'-s|=\Delta l$。
镜组透镜公式 $1/f=1/s+1/s'$ 给
$f=ss'/(s+s')$。
(a) $ss'=(l^2-\Delta l^2)/4$，所以
\[
\boxed{f=\frac{l^2-\Delta l^2}{4l}
=\frac{90^2-30^2}{4(90)}=20\,\mathrm{cm}}.
\]
(b) 两位置横向放大率模为
$\beta=s'/s$ 与其倒数；
比值 $\eta=\beta^2$，从而
$\beta=\sqrt\eta=2$。
取 $s=l/(1+\beta)=30$cm、
$s'=60$cm，得
\[
\boxed{f=\frac{ss'}l=20\,\mathrm{cm}}.
\]
\excheck 只在两个共轭实像位置存在时可用该方法，即 $l\ge4f$。

\opticsproblem{37}{193}
固定的物体与屏幕之间存在两种薄会聚透镜位置，
所得清晰像的横向尺寸分别为
$h_1'=2.0$mm 和 $h_2'=4.5$mm，
求物体本身横向尺寸 $h$。
\exsolution
两共轭透镜位置的横向放大率分别为
$\beta=s'/s$ 与
$\beta'=s/s'$，它们的乘积为1。
因此像大小之积
$h_1'h_2'=(\beta h)(h/\beta)=h^2$，故
\[
\boxed{h=\sqrt{h_1'h_2'}
=\sqrt{2.0\times4.5}\,\mathrm{mm}
=3.0\,\mathrm{mm}}.
\]
\excheck 物体大小是两个像尺寸的\emph{几何平均}，并非算术平均。

\opticsproblem{38}{193}
远处物体亮度 $L=260\,\mathrm{cd/m^2}$，
薄会聚摄影物镜孔径比
$D:f=1:3.5$；
透镜损耗光通量比例 $\alpha=0.10$。
求照片底片上成像照度。
\exsolution
对于远物，像距约为 $f$，透镜通光孔径在像点张的立体角近似
$\Omega\simeq\pi(D/2)^2/f^2$。
在近轴条件下像照度为
$E=(1-\alpha)L\Omega$，故
\[
\boxed{E=(1-\alpha)\frac{\pi LD^2}{4f^2}
=\frac{(0.9)\pi(260)}{4(3.5)^2}
\simeq15\,\mathrm{lx}}.
\]
\excheck 给的是摄影镜头的\emph{直径比} D/f，不是面积比；因此在照度公式中要平方。

\opticsproblem{39}{193}
薄会聚透镜成清晰实像，改变透镜通光口径 $D$，
讨论所成像的亮度：
(a) 眼睛直接观察真实空间像；
(b) 将像投到按朗伯定律漫反射的白屏上再观察。
\exsolution
(a) 在物镜光束完全进入眼瞳且未发生额外遮挡的条件下，
无耗损理想光学系统中辐亮度/亮度沿光线传播守恒。
改变口径主要改变通过的总光通量及视场，
并不改变被直接观察到的扩展实像表面亮度：
\[
\boxed{L_{\rm direct}\ \text{与}\ D\ \text{无关}}.
\]
(b) 屏上的照度由透镜收集的入射光通量决定，
通光面积 $\propto D^2$，
若其余条件固定，则
$E_{\rm screen}\propto D^2$。
朗伯白屏的反射亮度
$L_{\rm screen}=\rho E_{\rm screen}/\pi$，所以
\[
\boxed{L_{\rm screen}\propto D^2}.
\]
\excheck 两种观察方式须分别考虑“空间像的亮度不变量”与“漫反射屏上的照度”。

\opticsproblem{40}{193}
两只薄对称透镜紧密靠合，一只为双凸透镜
$n_1=1.70$，另一只为双凹透镜
$n_2=1.51$，两者表面曲率半径模均为
$R=10$cm。
将系统整体浸入水中，求像方焦距。
\exsolution
取水折射率 $n_w\simeq1.33$，
两透镜在水中的光焦度分别为
$\Phi_1=2(n_1-n_w)/R$，
$\Phi_2=-2(n_2-n_w)/R$。
紧贴时总光焦度
\[
\Phi=\Phi_1+\Phi_2=
\frac{2(n_1-n_2)}R
=\frac{2(1.70-1.51)}{0.10}
=3.8\,\mathrm D.
\]
介质中像方焦距 $f'=n_w/\Phi$，
因此
\[
\boxed{f'\simeq0.35\,\mathrm m=35\,\mathrm{cm}}.
\]
\excheck 水介质折射率在两透镜\emph{光焦度求和时}抵消，但求介质中的焦距还需乘 $n_w$。

\opticsproblem{41}{193}
取一只薄对称双凸玻璃透镜，
表面曲率半径模
$R=40$cm，将其中一面镀银，形成带折射面与反射面的凹面镜。
求等效焦距。
\exsolution
光依次在未镀银的凸面折射、
在银面反射、再从原折射面出射。
近轴薄系统的光焦度可按两次折射与一次反射的会聚功率相加，设玻璃折射率 $n\simeq1.50$：
\[
\Phi_{\rm eff}=
\frac{2(n-1)}R+\frac2R+
\frac{2(n-1)}R
=\frac{2(2n-1)}R.
\]
因而
\[
\boxed{f=\frac1{\Phi_{\rm eff}}
=\frac{R}{2(2n-1)}
=\frac{40\,\mathrm{cm}}{2(3-1)}
=10\,\mathrm{cm}}.
\]
\excheck 光两次穿过玻璃折射面，不能将其视为仅一面镀银的普通球面镜而用 $f=R/2$。

\opticsproblem{42}{193--194}
图5.9：三只共轴薄透镜依次排列，
光焦度为
$\Phi_1=+10.0$D、
$\Phi_2=-10.0$D、
$\Phi_3=+10.0$D；
相邻距离均为5.0cm。
求：(a) 从左平行射入的光束在最后一只透镜后汇聚的位置；
(b) 距第一透镜左侧多远有一点，与其像关于三透镜系统对称地位于两端？
\exsolution
使用光线列向量
$(y,u)^{\mathsf T}$，
自由传播和薄透镜矩阵分别为
$\mathsf P(d)=\begin{pmatrix}1&d\\0&1\end{pmatrix}$、
$\mathsf L(\Phi)=\begin{pmatrix}1&0\\-\Phi&1\end{pmatrix}$。
由左向右全系统
\[
\mathsf M=
\mathsf L(10)\mathsf P(0.05)
\mathsf L(-10)\mathsf P(0.05)\mathsf L(10)
=\begin{pmatrix}1/4&1/8\\-15/2&1/4\end{pmatrix}.
\]
(a) 平行光 $u_0=0$，
最后出射高度 $y'=y_0/4$、
斜率 $u'=-7.5y_0$，
向右传播到光轴需距离
\[
\boxed{s'=-y'/u'=\frac1{30}\,\mathrm m
\simeq3.3\,\mathrm{cm}}.
\]
(b) 若物点距第一透镜左侧 $\ell$，
像点距最后透镜右侧相同 $\ell$，
全传递矩阵
$\mathsf P(\ell)\mathsf M\mathsf P(\ell)$
的 (1,2) 元必须为零：
$-\tfrac{15}{2}\ell^2+\tfrac12\ell+\tfrac18=0$。
正根
\[
\boxed{\ell=\frac16\,\mathrm m\simeq17\,\mathrm{cm}}.
\]
\excheck 三透镜相隔各5cm，不能把总间距误当5cm；矩阵给出可重现的精确图示解。

\opticsproblem{43}{194}
伽利略望远镜角放大率为10，调焦到无穷远时筒长45cm。
求：(a) 物镜、目镜焦距；
(b) 观察50m远物体需将目镜移动多远。
\exsolution
伽利略望远镜物镜为正、目镜为负：
$\Gamma=f_{\rm obj}/|f_{\rm eye}|=10$，
调至无穷远筒长
$L=f_{\rm obj}-|f_{\rm eye}|=45$cm。
解得
\[
\boxed{f_{\rm obj}=50\,\mathrm{cm}},
\qquad
\boxed{f_{\rm eye}=-5.0\,\mathrm{cm}}.
\]
观察 $s=50$m 处物体，
物镜成像距
$s'=f_{\rm obj}s/(s-f_{\rm obj})$，
由原来0.500m变为
$0.50505$m，
因此
\[
\boxed{|\Delta x|=s'-f_{\rm obj}
\simeq0.00505\,\mathrm m=0.50\,\mathrm{cm}}.
\]
\excheck 远物变为有限距离时物镜中间像向右移动，需相应调整目镜的位置以恢复出射平行。

\opticsproblem{44}{194}
开普勒望远镜调至无穷远，物镜镜框直径为
$D$，经目镜成的镜框像直径为
$d$。求望远镜角放大率。
\exsolution
物镜所在平面经目镜的像称出瞳，
对于无穷远调整的望远镜，
光束的角放大率与光束口径缩小倍数互为对应：
$\Gamma=d\theta_{\rm out}/d\theta_{\rm in}
=f_{\rm obj}/f_{\rm eye}$。
从几何相似三角形得到出瞳直径
$d=D f_{\rm eye}/f_{\rm obj}$。
所以
\[
\boxed{\Gamma=\frac Dd}.
\]
\excheck 该结果需忽略遮挡及瞳孔截光，且把镜框的像正确视为出瞳。

\opticsproblem{45}{194}
望远镜使远方目标的光通量角集中度增大
$\eta=4.0\times10^4$ 倍。
目标经望远镜后的视角为
$\psi'=2.0^\circ$，
求目标自身视角 $\psi$。
\exsolution
在光损耗可忽略且面积源的小角度近似下，
角放大率为
$\Gamma=\sqrt\eta$。
因此
\[
\psi=\frac{\psi'}{\Gamma}
=\frac{2.0^\circ}{\sqrt{4.0\times10^4}}
=\frac{2^\circ}{200}
=0.01^\circ.
\]
\[
\boxed{\psi=0.60'}
\]
其中角分 $1^\circ=60'$。
\excheck 不能把能量强度倍数4万直接当角放大率；角尺度对应平方根关系。
